\documentclass[
 reprint,
 amsmath,amssymb,
 aps,
 prl,
 superscriptaddress,
 nofootinbib
]{revtex4-2}

\usepackage{graphicx}
\usepackage{bm}
\usepackage{xcolor}
\usepackage{placeins}
\usepackage[colorlinks=true,citecolor=blue!55!black,
linkcolor=blue!55!black,urlcolor=blue!55!black]{hyperref}

\newcommand{\Om}{\Omega}
\newcommand{\Res}{\operatorname*{Res}}
\newcommand{\cR}{\mathcal{R}}
\newcommand{\ii}{\mathrm{i}}
\newcommand{\crossed}{\text{crossed}}

\begin{document}

\title{Quasinormal-Mode Collider Spectroscopy}

\author{Mian Zhu}
\email{zhumian@scu.edu.cn}
\affiliation{College of Physics, Sichuan University, Chengdu 610065, China}

\date{\today}

\begin{abstract}
Black-hole spectroscopy uses the linear quasinormal-mode (QNM) spectrum to infer the geometry of a merger remnant, but nonlinear ringdown can also reveal how those modes interact. We introduce \emph{quasinormal-mode collider spectroscopy} (QCS), which organizes higher-order retarded responses by their analytic structure. A local quadratic interaction produces a composite sum-frequency pole. An exchanged QNM instead produces a pole in a pair-frequency channel whose residue factorizes into production and conversion vertices, directly paralleling particle exchange in cosmological collider correlators. We
validate the quadratic layer with the known Schwarzschild $220\times220\rightarrow440$ response and show that GW250114 already motivates a closure test using frequency, damping, phase, and amplitude together. For the exchange layer, we identify the near-extremal Kerr chain $210+220\rightarrow431\rightarrow651$, whose internal and output detunings are simultaneously small. QCS thereby promotes ringdown spectroscopy from measuring a mode spectrum to probing strong-gravity interactions.
\end{abstract}

\maketitle

\textit{Introduction.---}
The observation of gravitational waves (GWs) from binary-black-hole mergers opens a new avenue for testing gravity. Of particular interest is the post-merger ringdown, whose quasinormal-mode (QNM) spectrum encodes the properties of the remnant black hole \cite{Berti2009Quasinormal,London2014ModelingRingdown,Cardoso2016EventHorizon,Bhagwat2020Subdominant} and is now being probed with increasing precision by current and future GW detectors
\cite{LIGO2016Observation,LIGO2021GWTC2,LVK2023GWTC3,LVK2025GW250114,Punturo2010EinsteinTelescope,Luo2016TianQin,Hu2017Taiji,AmaroSeoane2017LISA,Reitze2019CosmicExplorer}. Standard black-hole spectroscopy treats each QNM as an independent linear oscillation and has been remarkably successful, for example in testing the Kerr nature of the remnant \cite{Isi2019NoHair} and identifying subtle overtones \cite{Giesler2019Overtones}.

This linear framework, however, neglects interactions among the modes. Such interactions can fix nontrivial amplitude ratios and relative phases that are sensitive to the interaction vertices of the theory, allowing nonlinear ringdown to probe the field equations directly. Recent analyses of GW250114
have reported evidence for nonlinear components of the QNM spectrum generated by the dominant quadrupolar ($\ell=m=2$) oscillation and its overtones \cite{LVKGW250114,YangShiHu2025,WangEtAl2026}. Nonlinear ringdown has therefore received increasing attention as a probe of gravity \cite{MitmanEtAl2023,CheungEtAl2024,BucciottiEtAl2024PRD,BucciottiEtAl2024JHEP,MaYang2024,KheraMaYang2025,BourgEtAl2025}.

In this Letter, we introduce quasinormal-mode collider spectroscopy (QCS), a framework for nonlinear black-hole spectroscopy. Our starting point is a structural analogy between black-hole ringdown and primordial cosmological fluctuations: both describe perturbations of a background spacetime that decompose into elementary modes at linear order, while nonlinearities couple these modes and generate higher-order signals. More precisely, the linear QNM spectrum appears as poles of the retarded Green function, which plays the same free-sector role as the cosmological two-point function; correspondingly, nonlinear ringdown responses are structurally analogous to the higher-point correlators that encode primordial non-Gaussianity.

Primordial non-Gaussianity has been studied extensively, providing systematic tools for extracting signatures of new physics from higher-point cosmological correlators. In particular, cosmological collider physics has shown how these correlators can reveal the properties and interactions of intermediate particles in the early universe  \cite{ArkaniHamedMaldacena2015,BaumannEtAl2021}. A central result of cosmological collider physics is that the exchange of an intermediate particle produces a characteristic singularity in a partial-momentum channel whose residue factorizes into its production and decay contributions. Analogously, we show that the exchange of an intermediate QNM produces a pole in a pair-frequency channel whose residue factorizes into the couplings governing its production and conversion. In this sense, the ringing black hole acts as a collider whose nonlinear response reveals the spectrum and interactions of intermediate QNMs, motivating the name QCS. As we shall show in this Letter, an exchanged QNM produces a pair-frequency pole whose residue factorizes into production and conversion couplings. We identify a pure-gravity near-extremal Kerr channel in which the intermediate and output modes are simultaneously near resonance, providing a concrete target for future calculations and gravitational-wave observations of high-spin remnants.


\begin{figure*}[t]
	\centering
	\includegraphics[width=\textwidth]{qcs.pdf}
	\caption{QCS hierarchy and its cosmological-collider analogue. The top row
	shows QNM responses and the bottom row cosmological correlators; the
	correspondence concerns diagrammatic roles, not equal observables or
	normalizations. (a) The linear QNM response and (d) primordial power
	spectrum represent the free sector. (b) A quadratic QNM response and (e) a
	contact bispectrum arise from one contact interaction, without an
	intermediate pole. (c) A QNM exchange pole and (f) an exchange trispectrum
	expose a factorizing subchannel singularity and its intermediate mode.}
	\label{fig:qcscc}
\end{figure*}


\textit{QCS principle.---}
We start with the linear QNM. Let $a$ denote an
asymptotic angular channel ( and also parity when needed), and let $h_a(\omega)$ be the corresponding strain multipole sourced by $J_a$. The retarded response near the pole is
\begin{equation}
	\label{eq:Green1}
	h_a^{(1)}(\omega)=G_a^R(\omega)J_a(\omega) ~,~
	G_a^R(\omega)\simeq\frac{Z_A}{\omega-\Om_A} ~,
\end{equation}
with $\Omega_A$ the complex QNM frequency satisfying $\operatorname{Im}\Om_A<0$ and $Z_A \equiv \Res_{\omega = \Om_A} G_a^R(\omega)$. Here, the retarded Green function $G_a^R$ generally contains a family of QNM poles indexed by the overtone number (n), so we introduce the uppercase letter $A = (a,n)$ to label a specific QNM pole, and the response for a given pole $A$ becomes 
\begin{equation}
	h_a^{(1)} (t) = \mathcal{A}_A (t_{\ast}) e^{-i \Omega_A (t - t_{\ast})} ~.
	\label{eq:QNMlinear}
\end{equation}

In cosmology, the analogous quantities of the retarded Green function $G_a^R$ are the Gaussian two-point function $\langle\zeta_{\bm{k}}\zeta_{\bm{k}'} \rangle =(2\pi)^3\delta^3(\bm{k}+\bm{k}')P_\zeta(k)$, which also encodes the full information of free sector. In addition, the cosmological power spectrum is an equal-time correlator of the mode function $u_k$, i.e., $P_\zeta(k)=|u_k(t_\star)|^2$, whereas $G_k^R(t,t') \propto \theta(t-t') \big[u_k(t)u_k^*(t') - \mathrm{c.c.} \big]$, see the first column of Fig. \ref{fig:qcscc}.

The first nonlinear QNM comes from the local quadratic interactions (Fig. \ref{fig:qcscc}(b)), whose response has the schematic form
\begin{align}
	h_a^{(2)}(\omega) & \nonumber = G_a^R(\omega) \sum_{bc} \int_{\omega_1,\omega_2}2\pi\delta(\omega-\omega_1-\omega_2) \\ & \times \Gamma^{(2)}_{a|bc}(\omega;\omega_1,\omega_2)h_b^{(1)}(\omega_1)h_c^{(1)}(\omega_2) ~.
	\label{eq:ha2}
\end{align}
This corresponds to a nonlinear vertex couples the three external responses with no intermediate propagating modes, thus analogues to  the contact diagram of three point correlation functions in cosmology (Fig.~\ref{fig:qcscc}(e)).

Using \eqref{eq:QNMlinear}, the delta function in \eqref{eq:ha2} immediately enforces the following composition
\begin{equation}
\Res_{\omega=\Om_B+\Om_C}h_a^{(2)}(\omega)
=G_a^R(\Om_B+\Om_C) \Gamma^{(2)}_{a|BC} \mathcal{A}_B \mathcal{A}_C ~,
\label{eq:composite}
\end{equation}
suggesting the factorization of non-linear response into parent excitation, nonlinear coupling, and the filtering of the nonlinear source by the remnant geometry. 

Following the same spirit, the next layer consists of exchanged QNMs. The intermediate mode must first be produced and then converted into the observed channel, so it appears generically in the third-order response with three input legs and one output leg (Fig.~\ref{fig:qcscc}(c)). After amputating the three
input and one output responses, the exchange contribution becomes
\begin{equation}
	\Gamma^{(3)}_{a|bcd} =\Gamma^{(2)}_{a|sd} G_s^R(\omega_b+\omega_c) 
	\Gamma^{(2)}_{s|bc} + \text{crossed channels} ~.
\end{equation}
Near a specific intermediate pole $S$, The full response therefore takes the form
\begin{equation}
\Gamma^{(3)}_{a|bcd}
=\Gamma^{(3),\mathrm{reg}}_{a|bcd}
+\frac{\Gamma^{(2)}_{a|Sd}\,Z_S\,\Gamma^{(2)}_{S|bc}}
{\omega_b+\omega_c-\Om_S}
+\crossed\ \text{channels} ~.
\label{eq:exchange}
\end{equation}
using \eqref{eq:Green1}. Here \(\Gamma^{(3),\mathrm{reg}}\) is analytic in the pair frequency, and an
uppercase index on a vertex leg denotes evaluation at the corresponding QNM
pole. Thus \(\Gamma^{(2)}_{S|bc}\) produces \(S\), while
\(\Gamma^{(2)}_{a|Sd}\) converts it into the output. The pole position gives
the intermediate complex frequency; its angular channel fixes angular momentum
and parity; and its factorized residue measures a product of quadratic
vertices. This nonanalytic structure cannot be reproduced by a local contact
interaction. At exact internal resonance, its time-domain limit is the secular
response \(t e^{-\ii\Om_S t}\).

Put parents \(B,C,D\) on shell, let \(A=(a,n_A)\) denote a candidate output
pole, and define \(\Om_T=\Om_B+\Om_C+\Om_D\),
\(\Delta_S=\Om_B+\Om_C-\Om_S\), and \(\Delta_A=\Om_T-\Om_A\). Restoring the
output response gives
\begin{equation}
\begin{split}
\mathcal A_a^{(3),\mathrm{ex}}
&=G_a^R(\Om_T)
\frac{\Gamma^{(2)}_{a|SD}Z_S\Gamma^{(2)}_{S|BC}}{\Delta_S}
\mathcal A_B\mathcal A_C\mathcal A_D\\
&\simeq\frac{Z_AZ_S\Gamma^{(2)}_{A|SD}\Gamma^{(2)}_{S|BC}}
{\Delta_A\Delta_S}\mathcal A_B\mathcal A_C\mathcal A_D.
\end{split}
\label{eq:double-pole}
\end{equation}
The two detunings determine the resonant kinematics, whereas the numerator
contains the interaction data. In a causal initial-value problem,
Eq.~(\ref{eq:double-pole}) is understood as a finite-difference response, not a
literal divergence. Exact double degeneracy gives
\(t^2e^{-\ii\Om_A t}/2\), up to residue and vertex factors (Supplemental
Material).

This factorizing pole is the direct analogue of an exchange singularity in a
cosmological four-point function~\cite{BaumannEtAl2021}, subject to three
essential differences: the ringdown observable is causal rather than in--out,
the QNM spectrum is complex because the black hole dissipates, and radial-mode
normalization must be removed by an asymptotic convention or by amputating the
response. Conserved QNM bilinear forms offer a natural route to computing the
Kerr residues in Eq.~(\ref{eq:exchange})~\cite{GreenEtAl2023}.

\begin{figure*}[t]
\includegraphics[width=\textwidth]{detuning.pdf}
\caption{\label{fig:qcs}
\textit{QNM collider kinematics.}
(a) The near-extremal chain \(210+220\rightarrow431\), followed by
\(431+220\rightarrow651\). (b) Complex detuning magnitudes from direct
continued-fraction root continuation. The shaded interval marks the double
near-resonance; circles and squares identify the internal and output minima,
respectively.}
\end{figure*}

\textit{Gravitational realizations.---}
We now anchor the two nonlinear layers in gravitational perturbation theory.
We use the shorthand $\ell mn$ for a mode with angular numbers $(\ell,m)$ and
overtone $n$; the fundamental quadrupolar mode is therefore 220. At quadratic
order, a known Schwarzschild self-coupling provides a direct benchmark. For
asymptotic strain in the spherical-harmonic convention,
\begin{equation}
\mathcal A_{440}^{(2)}
=0.15352\,e^{-0.06812\ii}\mathcal A_{220}^{\,2},
\label{eq:benchmark}
\end{equation}
in agreement with the published value
\(0.154e^{-0.068\ii}\)~\cite{BucciottiEtAl2024PRD}. This is the full dressed
strain response. Its decomposition into a Regge--Wheeler--Zerilli source
overlap and quadratic reconstruction terms depends on the choice of master
variable, but their sum in Eq.~(\ref{eq:benchmark}) does not. The benchmark
therefore shows that the dressed quadratic response is a calculable asymptotic
observable rather than a bare, convention-dependent vertex.

The exchange layer requires a genuine intermediate QNM. We therefore ask
whether the Kerr spectrum contains a pure-gravity channel for which both the
internal and output resonances are nearly satisfied. For spin weight
\(s=-2\), write \(\Om_{\ell mn}=\omega_{\ell mn}-\ii\gamma_{\ell mn}\). A scan
over fundamental parents with \(2\leq\ell\leq4\) and intermediate modes with
\(\ell_S\leq7\) and \(n_S\leq2\) selects the chain
\begin{equation}
210+220\rightarrow431,\qquad
431+220\rightarrow651.
\label{eq:chain}
\end{equation}
The two complex detunings are
\begin{equation}
\Delta_S=\Om_{210}+\Om_{220}-\Om_{431},\quad
\Delta_A=\Om_{210}+2\Om_{220}-\Om_{651}.
\label{eq:detunings}
\end{equation}
We first compute the frequencies with Leaver's continued-fraction method using
the public \texttt{qnm} implementation~\cite{Stein2019}, and then independently
continue and polish each root through the near-extremal region. At
\(\chi=0.99\),
\begin{align}
M\Delta_S&=0.012931+0.023120\ii,
&|M\Delta_S|&=0.026491, \nonumber\\
M\Delta_A&=0.001544-0.012505\ii,
&|M\Delta_A|&=0.012600 .
\label{eq:numbers}
\end{align}
The product is \(|M^2\Delta_S\Delta_A|=3.338\times10^{-4}\) and reaches
\(3.232\times10^{-4}\) near \(\chi=0.9890\). These values quantify only the
kinematic denominators of Eq.~(\ref{eq:double-pole}); by themselves they imply
neither a large response nor detectability.

Figure~\ref{fig:qcs} shows that both detunings remain small for
\(\chi\simeq0.988\)--\(0.992\). The continued-fraction residual stays below
\(1.4\times10^{-15}\) across the plotted interval. Direct root continuation is
essential in this regime because near-extremal Kerr spectra split into
zero-damping and damped branches~\cite{YangEtAl2013}.

The same chain also passes the angular selection rules. Azimuthal number is
conserved at both vertices, and the spherical-limit Wigner coefficients are
\[
\begin{pmatrix}4&2&2\\-3&1&2\end{pmatrix}
=-0.235702,\qquad
\begin{pmatrix}6&4&2\\-5&3&2\end{pmatrix}
=-0.226455 .
\]
Reflection parity is likewise allowed: the odd \(210\) and even \(220\) produce
an odd \(431\), which combines with a second even \(220\) to produce the odd
\(651\). Spheroidal mixing changes the numerical coefficients but does not
forbid the channel.

Equation~(\ref{eq:chain}) thus establishes resonant \emph{kinematics}, not a
signal amplitude. The physical response also depends on
\(Z_{651}Z_{431}\Gamma^{(2)}_{651|431,220}
\Gamma^{(2)}_{431|210,220}\), and on the
event-dependent excitation of the \(210\) parent. Near-extremal residue scaling
could either reinforce or offset the small detunings. Computing these
dynamical factors is the decisive next step, making this chain a concrete
target for Kerr perturbation theory rather than a claim of current
detectability.

\textit{Current observational reach.---}
Present observations do not reach the near-extremal exchange channel, but
GW250114 already illustrates how to test the quadratic layer of QCS. For a
component driven by two copies of the 220 mode, the complex-frequency sum rule
predicts
\(\Om_Q=2\Om_{220}\), hence
\(f_Q=2f_{220}\) and \(\tau_Q=\tau_{220}/2\). The LVK analysis of GW250114 gives
\(f_{220}=251.7^{+5.1}_{-5.0}\,\mathrm{Hz}\) and
\(\tau_{220}=4.09^{+0.42}_{-0.38}\,\mathrm{ms}\)~\cite{LVKGW250114}.
Propagating these one-dimensional marginal intervals gives
\(f_Q=503.4^{+10.2}_{-10.0}\,\mathrm{Hz}\) and
\(\tau_Q=2.05^{+0.21}_{-0.19}\,\mathrm{ms}\). The same analysis, using a
linear-QNM full-signal template, gives
\(f_{440}=503^{+130}_{-185}\,\mathrm{Hz}\) and
\(\tau_{440}=4.7^{+2.1}_{-2.7}\,\mathrm{ms}\). The real frequencies overlap
almost exactly, whereas the central damping times differ
(Fig.~\ref{fig:gw250114}). Real frequency alone therefore cannot distinguish a
linear 440 mode from a quadratic $220\times220$ response.

This comparison is a reach estimate, not an event reanalysis, because the
quoted intervals are model-conditioned marginals from separate hypotheses. A
decisive inference must fit the linear and quadratic interpretations jointly
and test the full closure relation: frequency, damping, phase, and amplitude.
The Schwarzschild response in Eq.~(\ref{eq:benchmark}) also cannot be
transferred directly to GW250114. Its remnant has moderate spin, so any
amplitude comparison requires a Kerr quadratic response evaluated at the
inferred spin and defined in the same asymptotic convention~\cite{MaYang2024}.

\begin{figure*}[t]
\centering
\includegraphics[width=\textwidth]{Gw.pdf}
\caption{\label{fig:gw250114}
\textit{Current reach of a quadratic QCS closure test with GW250114.}
Separate marginal 90\% intervals for the quadratic
\(220\times220\) prediction and the LVK linear-440 fit. Left: their frequency
intervals overlap. Right: their central damping times differ by
$2.65\,\mathrm{ms}$.
These one-dimensional marginals do not constitute a joint model comparison; a
quadratic identification additionally requires phase and amplitude closure.}
\end{figure*}

\textit{Discussion.---}
QCS extends black-hole spectroscopy from identifying a spectrum to measuring
interactions among its modes. At the quadratic layer, a composite pole defines
a dressed three-leg response coefficient. At the exchange layer, a
factorizing pair-frequency pole in the four-leg response identifies an
intermediate QNM: its complex position records frequency and lifetime, its two
vertices enforce angular selection rules, and its residue contains
strong-gravity interaction data.

Each layer has a distinct, falsifiable signature. A quadratic contribution
must satisfy the complex-frequency sum rule and exhibit the corresponding
composite pole. An exchange contribution must contain a
pair-frequency pole distinct from both the total composite pole and the output
resonance, and its residue must factorize as in Eq.~(\ref{eq:exchange}). These
criteria organize nonlinear waveforms in terms of invariant response data
rather than treating arbitrary raw ringdown amplitudes as universal.

The two gravitational examples anchor this hierarchy at complementary stages.
The known Schwarzschild result validates the dressed quadratic response.
The near-extremal chain in Eq.~(\ref{eq:chain}) supplies concrete exchange
kinematics, while leaving its dynamics to be determined: Kerr residues and
quadratic vertices will decide whether the double near-resonance yields a
measurable enhancement.

Current data reach the quadratic layer. For GW250114,
\(2\operatorname{Re}\Om_{220}\simeq\operatorname{Re}\Om_{440}\), so the real
frequency is degenerate, but damping together with phase and amplitude closure
can separate the two interpretations. Resolving a pair-frequency pole will
likely require either a population of remnants or a controlled
numerical-relativity scan across parent excitations. Additional fields would
introduce further internal poles and extend the same logic beyond the
gravitational QNM spectrum. In this way, higher-order ringdown becomes a
systematic probe of nonlinear black-hole dynamics.

\bibliography{ref}

\end{document}
